\section{Computational price-support diagnostics}\label{sec:computational}

The computational question is whether a physical optimum is supported by its
own marginal price. A useful diagnostic therefore reports the physical value
$D$, hull value $CH$, their nonnegative difference and own-price regret together.
We first strengthen the public lower bounds analytically, then summarize an
existing development screen. All cases in that screen are development cases;
two depot variants still represent one public timetable.

\subsection{An analytical baseline for both markets}

Use the public physical cases in Section~\ref{sec:public}. Market 0 has
$a_{0,t}=0.20$ throughout; market 1 has $a_{1,t}=0.18$ for the first 15 hours
and $0.22$ for the last 15. Both have $b=1/900$ in
$F_s(e)=\sum_t(a_{s,t}e_t+b e_t^2/2)$. This changes only supply cost, leaving
fleet feasibility fixed. As in the existing exact public certificates,
arithmetic uses the stored input coefficients as rationals; the decimal
parameters here describe those inputs to their displayed precision.

Let $E_{\min}$ be the two-bus total-energy floor obtained from the exact flow
relaxation, and let $E_{\rm svc}$ be mandatory service energy. With full
replenishment, efficiency one and zero monetary travel cost,
$E_{\min}=(L_{\rm flow}-2f)/a_0$. This gives about 1024.62 and 1071.97 kWh
for depots 15 and 16, respectively; $E_{\rm svc}$ is about 889.12 kWh in
both. The one-bus obstruction is unchanged. For a uniform price $p>0$,
every two-bus plan has $c+p\sum_t e_t\geq2f+pE_{\min}$, whereas every
plan using at least three buses has
$c+p\sum_t e_t\geq3f+pE_{\rm svc}$. Hence, whenever
\begin{equation}\label{eq:public-cardinality-check}
3f+pE_{\rm svc}\ \geq\ 2f+pE_{\min},
\end{equation}
the lower bound $2f+pE_{\min}$ holds for every physical plan and its convex
hull. For $p\geq\max_t a_{s,t}$, Fenchel's inequality then gives
\begin{equation}\label{eq:public-energy-floor}
CH_s\ \geq\ 2f+pE_{\min}
       -\sum_{t=1}^{30}\frac{(p-a_{s,t})^2}{2b}.
\end{equation}
The maximizing uniform test price is
$p_s^*=\bar a_s+bE_{\min}/30$, where $\bar a_s=\sum_t a_{s,t}/30$.
For both markets it exceeds 0.22 in each depot case, and the left minus
right side of~\eqref{eq:public-cardinality-check} is strictly positive:
about 67.76 at depot 15 and 56.17 at depot 16. Thus no unexamined fleet
cardinality invalidates the bound.

\begin{table}[t]
\centering\small
\caption{Analytical hull lower bounds for the ideal public model, rounded
down to two decimals. The last column is a numerical lower bound from the
separate later solver-baseline comparison, included only to expose the
missing analytical baseline. No runtime comparison or exact native-model
certificate is implied.}\label{tab:analytic-baseline}
\begin{tabular}{lrrr}
\toprule
Depot & Analytical, market 0 & Analytical, market 1 & Later computed, market 1\\
\midrule
15 & 424.36 & 418.96 & 405.83\\
16 & 435.67 & 430.27 & 420.45\\
\bottomrule
\end{tabular}
\end{table}

Table~\ref{tab:analytic-baseline} recovers the existing depot-15 original-market
certificate and extends it to the other depot and market. In the changed
market, the exact floors exceed the later numerical lower endpoints by
about 13.13 and 9.82 cost units. These inexpensive bounds should be included
before assessing what iterative optimization adds. The hypothetical higher
floor obtained by substituting a native flat optimum would retain native
model and tolerance qualifications; it is not used as an exact certificate.

\subsection{Physical cost, hull cost and incumbent regret}

The initial screen used the same two public cases and two small synthetic
cases: the two-service cyclic construction and a three-service case with
multiple depot visits. Each was evaluated in an original and a changed
market. Table~\ref{tab:joint-support} shows all four changed-market cases,
using the screen's physical planner, cold hull calculation and response at
the returned planner incumbent's own price, combined after the run. Public planner and hull stages had nominal 180-second wall allowances and
160-second native phases; responses had 60 seconds. Synthetic planner and hull
stages had 60-second wall allowances with 45-second native phases, and responses
had 30 seconds. Complete wrapper allowances were recorded separately. Call
and arithmetic limits also restricted the calculations. In the four-hour
synthetic target markets, the cyclic case uses $a_1=(0,3.8,0,0.2)$ and
$b=(0,0.2,0,0.2)$; the multivisit case uses $a_1=(0.18,0.18,0.22,0.22)$
and $b=(0.2,0.2,0.2,0.2)$. The archive fixes their physical inputs.

For the specific planner incumbent $\hat x$, let
$B_{\hat x}=c(\hat x)+p_{\hat x}^{\top}e(\hat x)$. A response enclosure
$V(p_{\hat x})\in[L_V,U_V]$ gives
$r(\hat x)\in[B_{\hat x}-U_V,B_{\hat x}-L_V]$.
The stored response price matches the gradient computed for that incumbent.
We also use the response lower bound to improve the hull lower bound through
$L_V-F_s^*(p_{\hat x})$, and any physical upper witness as a hull upper witness.
The physical lower bound cannot be weaker than the hull lower bound. The table
combines these same-screen inequalities before computing the gap interval;
it claims no reduction in the original algorithm's time.
These are native numerical results: they retain solver and physical-replay
tolerances and use rounded prices. Exact processing of their stored values
does not make them ideal-model proofs.

% Generated by research-20260928/review-response-v07/build_joint_table.py
\begin{table}[t]
\centering\small
\setlength{\tabcolsep}{4pt}
\caption{Changed-market support diagnostics from the initial development screen,
with same-screen native bounds combined after the run. All endpoints are rounded
outward. Regret belongs to the returned planner incumbent $\hat x$ in that row;
public planner optima are unknown. Analytical ideal-model energy floors are
reported separately in Table~\ref{tab:analytic-baseline}.}\label{tab:joint-support}
\begin{tabular}{lrrrr}
\toprule
Case & Physical $D$ & Hull $CH$ & Gap $D-CH$ & Regret $r(\hat x)$\\
\midrule
Two-service cyclic & $[98.99,99.01]$ & $[97.98,97.99]$ & $[1.01,1.02]$ & $[8.99,9.01]$\\
Three-service visits & $[84.71,84.73]$ & $[84.62,84.73]$ & $[0.00,0.10]$ & $[0.09,0.10]$\\
Public depot 15 & $[339.35,517.33]$ & $[339.35,517.33]$ & $[0.00,177.98]$ & $[190.43,235.09]$\\
Public depot 16 & $[377.65,547.88]$ & $[377.65,547.88]$ & $[0.00,170.22]$ & $[215.95,269.16]$\\
\bottomrule
\end{tabular}
\end{table}


The public rows establish sizeable incentives to change the \emph{returned}
schedules under their fixed own prices. They do not establish regret of an
optimal public fleet: the planner bounds remain open, so these schedules
may simply be costly incumbents. The gap intervals include zero and do not
certify that the public timetable lacks a supporting physical plan.

The multivisit response supplies a hull lower bound of 84.62, reducing the
same-screen planning-gap upper bound to 0.10 without another optimization.
This is a consequence of combining already available certificates, not a
measured benefit of reuse.

Postprocessing the same public planner uppers with the exact energy floors
from Table~\ref{tab:analytic-baseline} reduces the changed-market gap caps to
98.37 at depot 15 and 117.61 at depot 16. These are mixed-evidence bounds,
conditional on the native upper witnesses representing feasible plans of the
ideal model; the lower floors themselves are exact. They are not timed results
of the original algorithm. The original-market exact depot-15 certificate
in~\eqref{eq:public-nonlinear} remains separate and stronger for that market.

\subsection{What the computations establish}

\begin{table}[htbp]
\centering\small
\caption{Evidence supporting the current conclusions.}\label{tab:evidence-map}
\begin{tabularx}{\linewidth}{>{\raggedright\arraybackslash}p{0.27\linewidth}X}
\toprule
Finding & Scope and remaining question\\
\midrule
Exact support obstruction & A fully replenished synthetic fleet and a specified
robust variant have a positive gap; replication separates the planning gap
from whole-operator regret.\\
Analytical public bounds & Flow, cardinality and energy inequalities give exact
ideal-model lower bounds in two markets. Both depot variants share one timetable.\\
Incumbent public incentives & Native responses improve the bill of each named
screen incumbent at its own price. Regret at a planner optimum and a positive
public planning gap remain unresolved.\\
Repeated-solve acceleration & Development comparisons are preserved in the
supplement. Cap dependence, software failures and unconfirmed start acceptance
prevent a general speedup claim. Retrieval and learned proposals remain untested.\\
\bottomrule
\end{tabularx}
\end{table}

The computation presently demonstrates useful bounds and incumbent incentives,
with the limits in Table~\ref{tab:evidence-map}. It does not measure the
prevalence of price-support failure across transport networks. Further work
must improve physical-pricing bounds, compare methods at common quality and
scale to independent timetables before making general computational claims.
